% Study 2, the achieved npm panel, from the committed panel files and the status file. Input from main.tex
% in the extension; every number traces to sources.md and is asserted by check_numbers.py. Macros are
% provided in case the file is compiled before main.tex defines them.
\providecommand{\qs}{QScore9}

\section{Study 2: the achieved npm panel}
\label{sec:s2panel}

The packages domain is npm. Its panel was frozen on 2026-10-15 01:57 UTC and verified at 02:00 UTC, inside
the reset readiness deadline of 2026-10-20 23:59 UTC, so the registered PyPI substitution was not invoked.
Everything in this section is computed from the context window alone, 2025-09-01 to 2026-06-02; as of
2026-10-15, a statement to be re-checked on the posting day against the same status file, no target value
of 2026-06-03 to 2026-08-31 had been opened.

\paragraph{Inventory and candidate walk.} The frozen inventory is the registry's replication feed read in
4,493 pages of at most 1,000 names between 2026-10-14 00:23 and 02:04 UTC: 4,486,927 package names, scopes
included. The walk visited the first 12,943 names in candidate-hash order and requested each one's registry
document: 10,000 were eligible, created before 2025-09-01; 2,918 were ineligible on their creation time or
frame; and 25 were excluded because the registry answered HTTP 404 for the name. The walk stopped when the
pool was full. Because it did not exhaust the inventory, the size of the eligible frame is estimated from
the visited prefix, as the inventory size times the eligible share of the visited names: about 3,466,682
packages, with exact bounds of 10,000 and 4,483,984. The candidate-pool inclusion probability of an eligible
package is correspondingly about 0.0029, with bounds 0.0022 and 1. These are estimates from a hash-ordered
prefix, not counts.

\paragraph{The context pull and the strata.} The 10,000 candidates' daily download counts from 2025-09-01 to
2026-06-02 were requested from the download-counts API; 9,998 answered completely and 2 were excluded
because the API answered HTTP 404 for the package (\texttt{@zitterorg/modi-quaerat-voluptas},
\texttt{walk-logger}), the only context-stage exclusions. Of the 12,762 attempts to that host, 9,998 returned
the range, 40 returned 404, 2,708 returned 429, 2 returned 504 and 14 timed out; every refusal was retried
later from the saved captures without excluding a candidate. On the stratum window of 2026-03-05 to
2026-06-02 the 9,998 complete candidates fell 2,431 in the dense stratum (zero fraction below 0.5), 7,558 in
the sparse stratum (0.5 to below 0.9) and 9 in the very sparse stratum (0.9 and above). The achieved panel
is therefore 209 series: 100 dense, 100 sparse and 9 very sparse, the first in selection-hash order within
each stratum, with 91 very-sparse slots unfilled (Table~\ref{tab:s2panel}). The 209 names were created
between 2011-08-26 and 2025-08-08; 80 of them are scoped names.

\begin{table}[t]
\centering\small
\caption{The achieved npm panel by context-zero-fraction stratum, computed on 2026-03-05 to 2026-06-02:
the complete candidates among the 9,998 of the frozen pool, the series taken and the unfilled slots. The
stratum bounds are the registered $[0, 0.5)$, $[0.5, 0.9)$ and $[0.9, 1]$.}
\label{tab:s2panel}
\begin{tabular}{llrrr}
\toprule
Stratum & Zero fraction & Complete candidates & Taken & Unfilled \\
\midrule
dense & $[0, 0.5)$ & 2,431 & 100 & 0 \\
sparse & $[0.5, 0.9)$ & 7,558 & 100 & 0 \\
very sparse & $[0.9, 1]$ & 9 & 9 & 91 \\
\bottomrule
\end{tabular}
\end{table}

\paragraph{Why the third stratum is short.} The registered rule fixes the candidate pool at the first
10,000 eligible names in hash order and forbids a larger pool, a transfer of unused slots and any
alternative filter, so the stratum holds every complete candidate that qualified and no more. Only 9 of
9,998 complete candidates had zero downloads on at least 90 percent of the 90 stratum days. A reading, not a
measured fact: npm counts every tarball request, including mirrors and automated clients, so a published
package is rarely without a download for most of a quarter, and the very sparse stratum that is natural for
GitHub stars is close to empty for npm downloads. The stratum is reported as achieved; the confirmatory
family pools all 209 series as registered, strata being secondary descriptors, and any per-stratum table
shows $n = 9$ for the very sparse stratum.

\paragraph{The pull's budget and the API's limits.} The approval was 36,000 requests and 13 hours,
cumulative over the inventory, the walk and the context pull. The pull used 30,198 HTTP attempts: 12,943 to
the registry host, 12,762 to the download-counts host and 4,493 to the replication host, over 52
invocations totalling 45,508 seconds (12 hours and 38 minutes), from 2026-10-14 00:23 UTC to 2026-10-15 01:57
UTC. Attempts were spaced at least one second apart, and at least three seconds from 2026-10-14 22:51 UTC,
when the download API was answering about 20 requests a minute. The download API returned 429 to 2,708
attempts; three in a row for one key stopped the invocation, which resumed later from the saved captures.
From 2026-10-14 08:44 UTC the download API timed out or answered 504 until at least 10:25 UTC and answered
again at 21:37 UTC; each such stop wrote nothing and excluded no candidate. No documented request quota for
the download-counts API was found in the governing documents, and 30,198 attempts in about 26 hours is far
below the one stated line, five million requests a month; the redistribution review records the 429 answers
as the pull repeatedly reaching npm's throttle and honouring it.

\paragraph{One series at 7.6 million downloads a day.} The dense stratum holds the package
\texttt{aggregate-error}, whose last 28 context days before the first cutoff average 7,637,136 downloads a
day; the next largest last-28 mean in the panel is about 13,000. Equal series weights, as registered, mean
that the mean paired difference of the unscaled \qs{} over 209 series is that one series' difference divided
by 209 plus a remainder from the other 208, whose losses are of a different order; the series-bootstrap
interval inherits the same dependence. The median difference, the win count and the signed-rank test are
rank-based and unaffected; the scaled \qs{}, which divides each window's loss by the context mean, is
secondary. The paper states this when it reports the mean, and this is a property of the registered design,
not a defect found in it. The same series is the case behind the availability rule of 2026-10-15
(Section~\ref{sec:s2design}): no discrete CRPS can be certified for it within the registered work limit, so
its CRPS is recorded as unavailable and counted in the denominators.

\paragraph{Data statement.} The daily counts of the 209 series were obtained from the npm download-counts
API, \url{https://api.npmjs.org/downloads/range}, between 2026-10-14 and 2026-10-15 UTC, for 2025-09-01 to
2026-06-02; the targets for 2026-06-03 to 2026-08-31 are to be obtained under a separate target-access
approval. Package names and creation timestamps come from the npm public registry,
\url{https://registry.npmjs.org}, and from its replication feed, \url{https://replicate.npmjs.com}. A
download is a tarball request as npm counts it, including mirrors, build servers and automated clients, and
is a directional indicator rather than a count of users. npm, Inc.\ states no licence for the download
counts, and its Open-Source Terms of 2022-03-10, read as governing the public registry and its public APIs,
contain no restriction on their redistribution; replication of registry data through the public
APIs is expressly allowed. Names are as the registry served them at acquisition and may since have been
removed. The pull reached npm's rate limit 2,708 times and honoured it. The project's compilation of the
selection, the derived fields and the documentation is released under CC BY 4.0; no licence is asserted over
the counts as such. Maintainer names, e-mail addresses, descriptions and raw registry documents are not
published. The manifest and selection files are in the repository; the aggregated daily counts are
published with the paper. npm is a registered trademark of npm, Inc.; no npm logo is used.
